% =====================================================================
%  2bat-ud4-1.tex  —  SESSIÓ 1: El ritme de canvi i el vèrtex de 1r (connexió)
%  (sense preàmbul: s'inclou des de main.tex)
% =====================================================================
\horatitol{Sessió 1 — El ritme de canvi i el vèrtex de 1r}%
{Recuperar el vèrtex de les paràboles de 1r per trobar un punt òptim i obrir la pregunta de la unitat: i si la funció no és una paràbola?}

\subsection{Idea clau}

Obrim la darrera unitat d'\textbf{anàlisi} del curs. Connectem amb una cosa que ja
sabeu fer: a 1r, amb les paràboles, trobàveu el \textbf{vèrtex} per localitzar un punt
òptim (per exemple, el preu que maximitzava el benefici en una venda solidària). Avui
recuperem aquella idea i ens fem la pregunta que mou tota la unitat: \textbf{i si la
funció no és una paràbola? Com trobaríem el punt òptim de qualsevol funció?}

\begin{idea}
\textbf{Del vèrtex concret a la idea de ritme de canvi.} En el punt més alt (màxim) o
més baix (mínim) d'una corba, la gràfica \textbf{deixa de pujar i comença a baixar}
(o a l'inrevés). Just en aquell punt, per un instant, \textbf{no creix ni decreix}.
\begin{itemize}[leftmargin=1.4em, itemsep=2pt]
  \item Estudiar \emph{com creix o decreix} una funció a cada punt és estudiar el seu
        \textbf{ritme de canvi}.
  \item L'eina que mesura aquest ritme de canvi es diu \textbf{derivada}, i serà la
        protagonista de la unitat.
  \item El vèrtex de 1r serà, més endavant, només un \textbf{cas particular} d'un
        mètode general que servirà per a qualsevol funció.
\end{itemize}
\end{idea}

\subsection{Exemples}

\begin{exemple}[el vèrtex que ja sabíeu trobar: la venda solidària]
A 1r, el benefici d'una venda de samarretes solidàries segons el preu $x$ (en euros)
es modelitzava amb una paràbola, per exemple $B(x)=-2x^2+24x-40$. El benefici màxim
era al \textbf{vèrtex}, que es trobava amb la fórmula $x_v=\dfrac{-b}{2a}$:
\[
x_v=\frac{-24}{2\cdot(-2)}=\frac{-24}{-4}=6.
\]
El preu òptim era $6\eur$, i el benefici màxim $B(6)=-2\cdot 36+144-40=32\eur$. Aquell
vèrtex era el \textbf{punt òptim} del problema.
\end{exemple}

\begin{exemple}[i si no és una paràbola?]
La fórmula $x_v=\dfrac{-b}{2a}$ \textbf{només} funciona amb paràboles. Però molts
fenòmens socials no ho són:
\begin{itemize}[leftmargin=1.4em, itemsep=2pt]
  \item el cost per usuari d'un servei que primer baixa i després puja;
  \item els ingressos d'una activitat que pugen, arriben a un sostre i baixen;
  \item l'eficàcia d'un programa segons la seva durada.
\end{itemize}
Per a aquests casos necessitem una eina \textbf{general}: una que, mirant el
\textbf{ritme de canvi} a cada punt, ens digui on la funció deixa de créixer i comença
a decréixer. Aquesta eina és la \textbf{derivada}.
\end{exemple}

\subsection{Exercicis resolts}

\begin{exresolt}[recuperar el vèrtex d'una paràbola]
Una activitat solidària té un benefici $B(x)=-x^2+10x-9$ (en euros) segons el preu
d'entrada $x$ (en euros). Troba el preu que dóna el benefici màxim i quin és aquest
benefici.
\solucio
És una paràbola amb $a=-1$ i $b=10$. El vèrtex (màxim, perquè $a<0$) és a
\[
x_v=\frac{-b}{2a}=\frac{-10}{2\cdot(-1)}=\frac{-10}{-2}=5.
\]
El benefici màxim és $B(5)=-25+50-9=16\eur$. En el punt òptim, la gràfica deixa de
pujar i comença a baixar: el seu \textbf{ritme de canvi} és nul.
\resposta{Preu òptim $5\eur$; benefici màxim $16\eur$.}
\end{exresolt}

\begin{exresolt}[descriure el ritme de canvi amb paraules]
Una despesa acumulada d'un projecte social va, mes a mes, així: al principi puja
\emph{de pressa}, després puja \emph{a poc a poc}, i els últims mesos \emph{s'atura}
(no canvia). Descriu el ritme de canvi en cada tram i digues on és més gran.
\solucio
El \textbf{ritme de canvi} és la velocitat amb què canvia la despesa:
\begin{itemize}[leftmargin=1.4em, itemsep=2pt]
  \item Al principi, puja de pressa: el ritme de canvi és \textbf{gran i positiu}.
  \item Després, puja a poc a poc: el ritme és \textbf{positiu però petit}.
  \item Al final, s'atura: el ritme de canvi és \textbf{prop de zero}.
\end{itemize}
El ritme és més gran al principi, quan la despesa creix més de pressa.
\resposta{Ritme gran al principi (puja de pressa), petit després, gairebé nul al final
(s'atura).}
\end{exresolt}

\subsection{Exercicis per fer}

\begin{exercicis}
\begin{exlist}
  \item Una venda benèfica té un benefici $B(x)=-x^2+8x-7$ (en euros) segons el preu
        $x$. Troba el preu que el fa màxim i el benefici màxim, amb la fórmula del
        vèrtex de 1r.
  \item Per a la paràbola $f(x)=-2x^2+12x$, troba el vèrtex i digues si és un màxim o
        un mínim. (Pista: mira el signe de $a$.)
  \item Una despesa puja molt de pressa els primers dies, després cada cop més a poc a
        poc i finalment s'estabilitza. Ordena aquests tres moments de \emph{més} a
        \emph{menys} ritme de canvi.
  \item \textbf{(Interpretació)} Què vol dir, socialment, que l'atur d'un municipi
        «creixi cada cop més de pressa»? I què voldria dir que «creixés cada cop més a
        poc a poc»?
  \item \textbf{(Raonament)} Per què la fórmula $x_v=\dfrac{-b}{2a}$ no ens serveix per
        trobar el punt òptim d'una funció que no sigui una paràbola? Quina informació
        ens caldria, en general, per saber on una corba deixa de pujar?
  \item \textbf{(Reflexió)} Escriu amb les teves paraules què vol dir «ritme de canvi»
        i posa un exemple propi de l'àmbit social on interessi trobar un \emph{punt
        òptim} (un cost mínim, un benefici màxim, un aforament ideal\ldots).
\end{exlist}
\end{exercicis}

% ---------------------------------------------------------------------
\fullsolucions{Sessió 1}
\begin{solucions}
\begin{sollist}
  \item $a=-1$, $b=8$: $x_v=\dfrac{-8}{2\cdot(-1)}=4$. Benefici màxim
        $B(4)=-16+32-7=9\eur$. Preu òptim $4\eur$.
  \item $a=-2$, $b=12$: $x_v=\dfrac{-12}{2\cdot(-2)}=3$; $f(3)=-18+36=18$. Vèrtex
        $(3,18)$; com que $a<0$, és un \textbf{màxim}.
  \item De més a menys ritme: (1) els primers dies (puja molt de pressa), (2) després
        (puja més a poc a poc), (3) al final (s'estabilitza, ritme gairebé nul).
  \item Resposta oberta. Idea: «creix cada cop més de pressa» vol dir que el ritme de
        canvi (la velocitat de creixement) augmenta —cada mes s'hi afegeix més atur que
        l'anterior—; «cada cop més a poc a poc» vol dir que el ritme disminueix, encara
        que l'atur encara pugi.
  \item Resposta oberta. Idea: la fórmula del vèrtex està deduïda específicament per a
        paràboles ($ax^2+bx+c$); per a una corba qualsevol no hi ha aquesta simetria.
        En general ens cal una eina que mesuri el ritme de canvi a cada punt i ens digui
        on s'anul·la (on la corba deixa de pujar): aquesta eina és la derivada.
  \item Resposta oberta.
\end{sollist}
\end{solucions}
